\subsection{Methodology of data scrambling}
\label{lab:DataScrambling_method}

 In order to check data behavior with blinding on LIV effect, we have studied the data scrambling. 
This uses the measurement of Z counting for each LB in data. However, we do scrambling the time stamp. 
Because LIV effect can be probed by looking consequenced data, therefore if we scramble time stamp of each LB, 
we can not probe the LIV effect. 

 Actual method is very simple, we scramble time stamp of LBs within each year. No same time stamp is used twice. 
Actual function to scramble time stamps is the function, ''sample'', in random library of python3.

By using this dataset, we can check real value from the measurement with the lucid luminosity and the Z counting. 
In order to check staticial behavior, 1000 (or 10000) set of scrabled data are generated by using different seed in same way.
Number of LBs for each year of data is listed in Table.~\ref{tab:data_LBs}. 
Those are about 90 k LBs\footnote{we combine data15 and data16 because data15 has only 22 k LBs}, 
therefore we could generate different scriabled dataset by using different seed.

As we have three channels, electron, muon and combined, we check these output. 
